\chapter{Credit and Debit Valuation Adjustments}\label{ch:rc:credit-and-debit-valuation-adjustments}

When the Basel Committee finalised its capital rules for counterparty risk in June 2011,
it gave the reason in one sentence: during the financial crisis, roughly two-thirds of the
losses attributed to \omterm{def:rc:counterparty-exposure:ccr}{counterparty credit risk} came from changes in the \omterm{def:rc:credit-and-debit-valuation-adjustments:cva}{credit valuation
adjustment}, and only about one-third from actual defaults. Banks had lost money not
because their counterparties failed, but because the market price of the chance that they
might had risen, and that price was in the fair value of every uncollateralised derivative.
This chapter prices the chance: the \omterm{def:rc:credit-and-debit-valuation-adjustments:cva}{credit valuation adjustment} that a bank deducts from the
value of its derivatives for the counterparty's default, the \omterm{def:rc:credit-and-debit-valuation-adjustments:dva}{debit valuation adjustment} it
adds for its own, how both are computed from chapter~17's exposures and chapter~13's hazard
curves, and how they are hedged, accounted for and treated in capital.

\section{CVA as the price of counterparty default}

\begin{definition}[Credit valuation adjustment]\label{def:rc:credit-and-debit-valuation-adjustments:cva}
The \emph{credit valuation adjustment}\index{credit valuation adjustment} (CVA) of a \omterm{def:rc:counterparty-exposure:netting}{netting
set} is the market value of the loss from the counterparty's default: the value of the trades
without counterparty risk minus their value with it. With the counterparty's default time $\tau_C$,
recovery $R_C$ and discounted exposure,
\begin{align*}
\mathrm{CVA} &= (1-R_C)\,\E\bigl[\mathbf 1_{\tau_C\le T}\,D(0,\tau_C)\max(V_{\tau_C}-C_{\tau_C},0)\bigr]\\
&\approx (1-R_C)\sum_k\mathrm{DEE}(t_k)\bigl(Q_C(t_{k-1})-Q_C(t_k)\bigr),
\end{align*}
where $\mathrm{DEE}(t) = \E[D(0,t)E_t]$ and the approximation assumes the exposure and the
default independent.
\end{definition}

The formula is a sum over default dates of three factors: the \omterm{def:rc:reduced-form-credit:pd}{loss given default}, the
probability that default falls in each period (from the counterparty's default-swap curve),
and the discounted \omterm{def:rc:counterparty-exposure:profiles}{expected exposure} then. It is the value of a contingent default swap
whose notional is the future exposure.

\begin{example}[The netting set's CVA]\label{ex:rc:credit-and-debit-valuation-adjustments:cva}
The counterparty of chapter~17's \omterm{def:rc:counterparty-exposure:netting}{netting set} is quoted at 80, 95, 110, 140, 160 and 175
basis points at one to ten years (illustrative BBB), recovery 40\%: its ten-year \omterm{def:rc:reduced-form-credit:pd}{probability
of default} is 26.4\%. With the exposures discounted path by path along the Hull--White
scenarios, the \omterm{def:rc:counterparty-exposure:netting}{netting set}'s CVA is USD~891\,342. Standalone, the swap's CVA is 288\,250 and
the cross-currency swap's 768\,334: netting saves 16\%. Over the ten-year default swap's \omterm{def:rc:reduced-form-credit:risky}{risky
annuity} of 7.45, the CVA is 12.0 basis points a year on USD~100 million
(\cref{fig:rc:credit-and-debit-valuation-adjustments:contrib}).
\end{example}

\begin{omfigure}
\begin{tikzpicture}
\begin{axis}[width=11.5cm,height=5.2cm,ybar,bar width=14pt,
  xlabel={year of default},ylabel={CVA contribution (USD million)},
  xtick={0.5,1.5,2.5,3.5,4.5,5.5,6.5,7.5,8.5,9.5},xticklabels={1,2,3,4,5,6,7,8,9,10},
  tick label style={font=\small},label style={font=\small},
  xmin=0,xmax=10,ymin=0,ymax=0.14,grid=major,grid style={gray!25},
  yticklabel style={/pgf/number format/fixed,/pgf/number format/precision=2},
  table/col sep=comma]
\addplot[fill=omDef!70,draw=omDef] table[x=year,y=cva]{figdata/rates-credit-risk/18-credit-and-debit-valuation-adjustments/contrib.csv};
\end{axis}
\end{tikzpicture}

\omcaption{The \omterm{def:rc:counterparty-exposure:netting}{netting set}'s CVA by year of the counterparty's default: the product of the
discounted \omterm{def:rc:counterparty-exposure:profiles}{expected exposure} (rising, then flat) and the \omterm{def:rc:reduced-form-credit:pd}{probability of default} in the year
(rising with the hazard curve, then discounted by survival). Years four to six contribute most.
Data: the chapter's tutorial.}\label{fig:rc:credit-and-debit-valuation-adjustments:contrib}
\end{omfigure}

\omcode{code/firm/cva/firm_cva.py}{42}{48}{CVA from a discounted expected exposure profile and a hazard curve, unilateral or first-to-default.}\label{lst:rc:credit-and-debit-valuation-adjustments:cva}

\section{DVA and the bilateral view}

\begin{definition}[Debit valuation adjustment]\label{def:rc:credit-and-debit-valuation-adjustments:dva}
The \emph{debit valuation adjustment}\index{debit valuation adjustment} (DVA) is the same
quantity seen from the counterparty: the market value of the gain to the bank from its own
default, when it would stop paying what it owes, computed from the discounted negative
exposure and the bank's own hazard curve and recovery.
\end{definition}

\begin{definition}[Bilateral CVA]\label{def:rc:credit-and-debit-valuation-adjustments:bcva}
The \emph{bilateral CVA}\index{bilateral CVA} is the adjustment $\mathrm{CVA}-\mathrm{DVA}$ when each is
computed on first-to-default: the counterparty's default costs the bank only if the bank
has not defaulted first (weight $Q_B(t)$ in the CVA sum), and symmetrically. It makes the two
parties agree on the price: one's CVA is the other's DVA.
\end{definition}

\begin{example}[Both sides]\label{ex:rc:credit-and-debit-valuation-adjustments:bcva}
The bank itself is quoted at 40 to 90 basis points (illustrative single-A). Its \omterm{def:rc:counterparty-exposure:netting}{netting set}
drifts in the counterparty's favour (chapter~17), so its DVA of USD~1\,055\,437 exceeds its CVA.
On first-to-default the CVA is 832\,416 and the DVA 896\,328: the bilateral adjustment is
$-\text{USD}~63\,912$, a small net \emph{benefit} to the bank. Under chapter~17's zero-threshold CSA
the CVA falls to USD~122\,340 and the DVA to 70\,490.
\end{example}

\begin{remark}[Can a bank monetise its own default?]\label{rem:rc:credit-and-debit-valuation-adjustments:dva}
DVA is real value to the bank's creditors, not to its shareholders: it is realised only by
defaulting. It cannot be hedged by buying protection on oneself; banks approximate the hedge
by selling protection on similar banks, which adds correlation risk. Its growth when the bank's
credit worsens is why regulators remove it from capital.
\end{remark}

\begin{definition}[Close-out amount]\label{def:rc:credit-and-debit-valuation-adjustments:closeout}
The \emph{close-out amount}\index{close-out amount} is the amount the surviving party
determines under the master agreement for the terminated \omterm{def:rc:counterparty-exposure:netting}{netting set}: broadly the cost of
replacing the transactions, which may include the replacement counterparty's view of the
survivor's own credit (a risky close-out) or not (a risk-free close-out). The choice changes the CVA formula: with a risky close-out, the
surviving party's DVA survives the counterparty's default.
\end{definition}

\section{Computing CVA}

\begin{method}[CVA in a simulation engine]\label{met:rc:credit-and-debit-valuation-adjustments:engine}
(1) Simulate exposures and deflators on common scenarios (chapter~17). (2) Compute the
discounted expected positive and negative exposures of each \omterm{def:rc:counterparty-exposure:netting}{netting set}. (3) Bootstrap each
counterparty's hazard curve from its default-swap quotes, or from a proxy curve for a name
without quotes (chapter~20). (4) Sum over the grid. (5) For \omterm{def:rc:counterparty-exposure:wwr}{wrong-way risk}, simulate the
hazard rate with the market factors and average the pathwise loss.
\end{method}

\begin{example}[Wrong-way CVA]\label{ex:rc:credit-and-debit-valuation-adjustments:wwr}
Tie the counterparty's hazard to the euro:
\[
\lambda_t = \lambda^0_t\,\bigl(X_t/F_X(0,t)\bigr)^{-b},
\]
normalised so
that the average survival curve is still the market's, as in chapter~17. The cross-currency
swap's CVA is USD~766\,532 without dependence ($b = 0$), 1\,108\,459 with $b = 2$ and 1\,521\,524 with
$b = 5$: \omterm{def:rc:counterparty-exposure:wwr}{wrong-way risk} doubles it (\cref{fig:rc:credit-and-debit-valuation-adjustments:wwr}).
With very strong dependence the CVA falls back a little: the counterparty then defaults early
on the paths where the euro falls first, before the exposure has grown.
\end{example}

\begin{omfigure}
\begin{tikzpicture}
\begin{axis}[width=11.5cm,height=5.2cm,
  xlabel={dependence $b$ of the hazard on the euro},ylabel={CVA (USD million)},
  tick label style={font=\small},label style={font=\small},
  xmin=0,xmax=10,ymin=0,ymax=1.8,grid=major,grid style={gray!25},
  table/col sep=comma]
\addplot[omThm,very thick,mark=*,mark size=1.5pt] table[x=b,y=cva]{figdata/rates-credit-risk/18-credit-and-debit-valuation-adjustments/wwr.csv};
\end{axis}
\end{tikzpicture}

\omcaption{CVA of the cross-currency swap as its counterparty's hazard depends more strongly
on the euro, with the average survival curve held at the market's. \omterm{def:rc:counterparty-exposure:wwr}{Wrong-way risk} doubles the
CVA; the curve turns down when defaults concentrate so early that the exposure is still small.
Data: the chapter's tutorial.}\label{fig:rc:credit-and-debit-valuation-adjustments:wwr}
\end{omfigure}

\section{Sensitivities and hedging}

CVA moves with everything the exposure depends on (rates, FX, volatilities) and with the
counterparty's credit. A CVA desk hedges the first set with the underlying instruments and
the second with default swaps on the counterparty, or on an index when the name does not
trade.

\begin{example}[Bucketed CS01]\label{ex:rc:credit-and-debit-valuation-adjustments:cs01}
Raising each of the counterparty's six quoted spreads by one basis point, with the curve
re-bootstrapped, changes the CVA by $-235$, $-165$, $-141$, $-22$, $+125$ and $+4\,693$ dollars
(\cref{fig:rc:credit-and-debit-valuation-adjustments:cs01}): a parallel rise costs
USD~4\,254 per basis point. Almost all of it sits on the ten-year quote; raising a short quote
with the long ones fixed lowers the later forward hazards, so short buckets are negative.
The hedge is about USD~5.7 million of ten-year protection (the \omterm{def:rc:reduced-form-credit:cs01}{CS01} over the default swap's
\omterm{def:rc:reduced-form-credit:risky}{risky annuity} and a basis point).
\end{example}

\begin{omfigure}
\begin{tikzpicture}
\begin{axis}[width=11.5cm,height=5.2cm,ybar,bar width=18pt,
  symbolic x coords={1y,2y,3y,5y,7y,10y},xtick=data,
  xlabel={counterparty par spread bumped},ylabel={change in CVA (USD per bp)},
  tick label style={font=\small},label style={font=\small},
  ymin=-1200,ymax=5600,grid=major,grid style={gray!25},
  nodes near coords,every node near coord/.append style={font=\footnotesize,/pgf/number format/fixed,/pgf/number format/precision=0},
  table/col sep=comma]
\addplot[fill=omThm!60,draw=omThm] table[x=tenor,y=cs01]{figdata/rates-credit-risk/18-credit-and-debit-valuation-adjustments/cs01.csv};
\end{axis}
\end{tikzpicture}

\omcaption{Bucketed \omterm{def:rc:reduced-form-credit:cs01}{CS01} of the \omterm{def:rc:counterparty-exposure:netting}{netting set}'s CVA. The ten-year quote carries almost all the
sensitivity; shorter quotes have small negative sensitivities through the bootstrap. Data: the
chapter's tutorial.}\label{fig:rc:credit-and-debit-valuation-adjustments:cs01}
\end{omfigure}

\section{Accounting and regulation}

\begin{definition}[XVA]\label{def:rc:credit-and-debit-valuation-adjustments:xva}
\emph{XVA}\index{XVA} is the family of valuation adjustments that turn a derivative's
riskless value into the price a bank charges or reports: CVA and DVA for default, and the
funding, margin and capital adjustments of chapter~19.
\end{definition}

Accounting standards value derivatives at fair value, an exit price: IFRS~13, issued in May
2011, defines it with the assumptions market participants would use, including about risk.
Counterparty credit is such a risk, so CVA is part of reported fair value, and so is the
bank's own credit, as DVA: the fair value of a liability reflects its non-performance risk, which
includes the entity's own credit risk (paragraph 42). JPMorgan's third quarter of 2011, for instance,
included a 1.9 billion dollar pretax DVA gain from the widening of its own credit spreads. Regulation takes a different view: CVA losses need capital (the
charge the Basel Committee introduced after the crisis, chapter~19), and DVA, which rises when
the bank weakens, is deducted from common equity.

\begin{dated}{2026-09}{CVA and DVA in capital}\label{dat:rc:credit-and-debit-valuation-adjustments:dva}
The Basel Committee stated in June 2011 that about two-thirds of the crisis losses from
counterparty risk were CVA losses. Basel~III requires banks to derecognise from common
equity tier~1 the unrealised gains and losses from changes in their own credit risk on fair
valued liabilities; its final rule of July 2012 extended this to all DVA on derivatives,
phased in from 20\% in 2014 to a full deduction from 1 January 2018.
\end{dated}

\begin{omfigure}
\begin{tikzpicture}
\begin{axis}[width=11.5cm,height=5.2cm,
  xlabel={widening of the bank's own spreads (bp)},ylabel={DVA (USD million)},
  tick label style={font=\small},label style={font=\small},
  xmin=0,xmax=100,ymin=0,ymax=2.2,grid=major,grid style={gray!25},
  table/col sep=comma]
\addplot[omDef,very thick,mark=*,mark size=1.5pt] table[x=shift,y=dva]{figdata/rates-credit-risk/18-credit-and-debit-valuation-adjustments/dva.csv};
\end{axis}
\end{tikzpicture}

\omcaption{DVA of the \omterm{def:rc:counterparty-exposure:netting}{netting set} as the bank's own default-swap spreads widen in parallel:
the worse the bank's credit, the larger the accounting gain. Data: the chapter's
tutorial.}\label{fig:rc:credit-and-debit-valuation-adjustments:dva}
\end{omfigure}

\section{Tutorial: CVA of a netting set}\label{tut:rc:credit-and-debit-valuation-adjustments:tutorial}

\begin{tutorial}
\textbf{Goal.} Compute the CVA and DVA of chapter~17's \omterm{def:rc:counterparty-exposure:netting}{netting set}, unilateral and
first-to-default, with and without a CSA, their credit sensitivities, and the effect of
\omterm{def:rc:counterparty-exposure:wwr}{wrong-way risk}. \textbf{End state:} the numbers of
\cref{ex:rc:credit-and-debit-valuation-adjustments:cva,ex:rc:credit-and-debit-valuation-adjustments:bcva,ex:rc:credit-and-debit-valuation-adjustments:wwr,ex:rc:credit-and-debit-valuation-adjustments:cs01}
and the four charts.
\begin{tutsteps}
\item \textbf{Inputs}: chapter~17's scenarios and values; \texttt{deflators(sc)};
\texttt{bootstrap} for both credit curves.
\item \textbf{Adjustments}: \texttt{adjustments()} and \texttt{adjustments(csa=True)}.
\item \textbf{Risk}: \texttt{cs01()}; \texttt{wrong\_way()}.
\item \textbf{DVA}: \texttt{dva\_shift()}; \texttt{fig\_rc\_cva.py} writes the charts.
\end{tutsteps}
\textbf{What to change next.} Replace the counterparty curve by a flat 175 basis points and
compare; compute the CVA with undiscounted EE times the zero-coupon bond, and measure the
error of ignoring the correlation of rates and exposure.
\end{tutorial}

\section{Build: the CVA engine}\label{bld:rc:credit-and-debit-valuation-adjustments:cva}

\begin{build}
\bfield{Purpose} The firm's CVA and DVA: fair-value adjustments of every \omterm{def:rc:counterparty-exposure:netting}{netting set}, their
credit sensitivities for the hedging desk, and the inputs of chapter~20's pricing.
\bfield{Interface} \texttt{deflators(sc)}; \texttt{discounted\_profiles(E, D)};
\texttt{cva(dee, times, cpty, recovery, own=None)}; \texttt{dva}; \texttt{cva\_pathwise(E, D,
times, hazards)}; \texttt{cs01\_buckets}; \texttt{running\_spread}.
\bfield{Rules} Exposure at period midpoints; independence unless a hazard is simulated;
recovery 40\% by default; curves from \texttt{firm\_cdscurve}.
\bfield{Acceptance tests} \texttt{code/firm/cva/tests/}: flat exposure and hazard give the
closed form; the pathwise engine with a deterministic hazard equals the profile engine;
first-to-default CVA is below unilateral; bucketed \omterm{def:rc:reduced-form-credit:cs01}{CS01} add up to the parallel \omterm{def:rc:reduced-form-credit:cs01}{CS01}.
\bfield{Stretch} Stochastic recovery; risky close-out; CVA Greeks by adjoint algorithmic
differentiation; incremental CVA (chapter~20).
\end{build}

\begin{omsources}
\item Basel Committee on Banking Supervision, press release of 1 June 2011 on the capital
treatment of bilateral counterparty credit risk.
\item Basel Committee on Banking Supervision, \emph{Basel~III}, revised June 2011,
paragraph~75; press release of 25 July 2012 on valuation adjustments to derivative
liabilities.
\item IFRS Foundation, IFRS~13 \emph{Fair Value Measurement}, 2011.
\end{omsources}

\section{Exercises}

\begin{exercise}[$\star$]\label{exo:rc:credit-and-debit-valuation-adjustments:1}
A \omterm{def:rc:counterparty-exposure:netting}{netting set} has a flat discounted EE of USD~2 million for five years; the counterparty's
hazard is a flat 2\% and its recovery 40\%. Give the CVA.
\end{exercise}

\begin{exercise}[$\star$]\label{exo:rc:credit-and-debit-valuation-adjustments:2}
Estimate the chapter's CVA as EPE times the ten-year spread times the \omterm{def:rc:reduced-form-credit:risky}{risky annuity}, and
compare.
\end{exercise}

\begin{exercise}[$\star$]\label{exo:rc:credit-and-debit-valuation-adjustments:3}
Why is the bilateral adjustment of \cref{ex:rc:credit-and-debit-valuation-adjustments:bcva}
negative, and what does the counterparty see?
\end{exercise}

\begin{exercise}[$\star\star$]\label{exo:rc:credit-and-debit-valuation-adjustments:4}
Why is the \omterm{def:rc:counterparty-exposure:netting}{netting set}'s CVA less than the sum of the standalone CVAs?
\end{exercise}

\begin{exercise}[$\star\star$]\label{exo:rc:credit-and-debit-valuation-adjustments:5}
How much ten-year protection hedges the CVA's credit sensitivity, and what does the hedge
leave open?
\end{exercise}

\begin{exercise}[$\star\star$]\label{exo:rc:credit-and-debit-valuation-adjustments:6}
Why does the CSA cut the CVA by 86\% but the DVA by 93\%?
\end{exercise}

\begin{exercise}[$\star\star\star$]\label{exo:rc:credit-and-debit-valuation-adjustments:7}
\emph{Coding.} Compute the wrong-way CVA for $b = 0$, 2 and 5, and explain the turn of
the curve at large $b$.
\end{exercise}

\begin{exercise}[$\star\star\star$]\label{exo:rc:credit-and-debit-valuation-adjustments:8}
\emph{Find the flaw.} ``Our DVA gain this quarter proves that our derivatives business is
more valuable than the market thought.''
\end{exercise}

\section{Problem: The Gain from Getting Riskier}

\begin{problem}[{Weekend problem --- own credit and the income statement}]\label{pb:rc:credit-and-debit-valuation-adjustments:1}
In a quarter of stress the bank's own default-swap spreads widen by 50 basis points at every
maturity while nothing else moves. Consider chapter~17's \omterm{def:rc:counterparty-exposure:netting}{netting set} and the bank's curve of
\cref{ex:rc:credit-and-debit-valuation-adjustments:bcva}.

\medskip\noindent\textbf{Part I --- The gain.}
\begin{enumerate}
\item Give the unilateral DVA before and after.
\item Give the gain, and say where it appears in the accounts.
\item Why does the DVA rise when the bank's credit worsens?
\item Scale the gain to a derivatives book 500 times larger.
\item Who is worse off in the same quarter?
\end{enumerate}

\medskip\noindent\textbf{Part II --- Capital.}
\begin{enumerate}[resume]
\item What does Basel~III do with the gain?
\item When was the deduction fully phased in?
\item Why do regulators treat DVA differently from CVA?
\item What would happen to the bank's common equity tier~1 ratio if the gain were counted?
\item How would investors read the gain?
\end{enumerate}

\medskip\noindent\textbf{Part III --- Hedging it.}
\begin{enumerate}[resume]
\item Why can the bank not buy protection on itself?
\item What proxy hedge could it use, and what risk does it add?
\item Should the bank hedge DVA at all?
\item How does a CSA change the DVA and the gain?
\item What does the counterparty see in its CVA?
\end{enumerate}

\medskip\noindent\textbf{Part IV --- Judgement.}
\begin{enumerate}[resume]
\item Is DVA value? For whom?
\item How should a trader's P\&L treat DVA?
\item Why is the bilateral adjustment the one both parties can agree on?
\item State the \emph{named result}: the DVA gain from a 50-basis-point widening on the
\omterm{def:rc:counterparty-exposure:netting}{netting set}.
\item In one sentence: why is DVA deducted from capital?
\end{enumerate}
\end{problem}

\section{Interview questions}

\begin{interviewq}[$\star$ \iqroles{trader, risk}]\label{iq:rc:credit-and-debit-valuation-adjustments:1}
Write the CVA formula and explain each term.
\end{interviewq}

\begin{interviewq}[$\star\star$ \iqroles{researcher}]\label{iq:rc:credit-and-debit-valuation-adjustments:2}
What is DVA, and why is it controversial?
\end{interviewq}

\begin{interviewq}[$\star\star$ \iqroles{trader}]\label{iq:rc:credit-and-debit-valuation-adjustments:3}
How would you hedge the CVA of a portfolio of uncollateralised swaps with a BBB corporate?
\end{interviewq}

\begin{interviewq}[$\star\star$ \iqroles{developer, researcher}]\label{iq:rc:credit-and-debit-valuation-adjustments:4}
How do you include \omterm{def:rc:counterparty-exposure:wwr}{wrong-way risk} in a CVA engine?
\end{interviewq}

\begin{interviewq}[$\star\star\star$ \iqroles{risk, bank}]\label{iq:rc:credit-and-debit-valuation-adjustments:5}
Why did CVA cause more losses than defaults in the crisis?
\end{interviewq}

\begin{interviewq}[$\star\star\star$ \iqroles{researcher}]\label{iq:rc:credit-and-debit-valuation-adjustments:6}
Risk-free versus risky close-out: what is the difference and why does it matter?
\end{interviewq}
